\section{Results}

\subsection{Data}

We analyze excess P\&I mortality from all US states over 12 flu
seasons (2010--2018 and 2023--2025, excluding the COVID-19 pandemic
years). Figure~\ref{fig:seasons} shows the raw weekly death counts
for four representative states, with the flu seasons used for
inference highlighted. The synchronized winter peaks across states are
clearly visible: all four time series rise and fall together each
year, precisely the regime our analysis predicts will confound
connectivity inference.

\begin{figure}[h]
\centering
\includegraphics[width=\textwidth]{../pix/seasons_highlighted.pdf}
\caption{Weekly P\&I deaths for four US states. Light traces show the
  full time series; darker segments indicate the 20-week flu seasons
  (starting each November) included in the analysis. COVID-19
  pandemic years (2020--2022) are excluded. A pre-season baseline is
  subtracted before model fitting.}
\label{fig:seasons}
\end{figure}

\subsection{Synchronization Inflates Data Requirements}

Figure~\ref{fig:crlb_ratio} shows the ratio of median CRLB
(synchronized / asynchronous) under the two IC regimes. When initial
conditions are similar (left), the ratio is large throughout the
parameter space: synchronization inflates the data requirement by
orders of magnitude. When initial conditions differ (right), the
ratio is much smaller---the two panels look qualitatively similar but
the color scales differ dramatically (note the colorbar values).
Different initial conditions break the symmetry that makes
synchronized and unsynchronized epidemics indistinguishable, reducing
the penalty imposed by synchronization.

\begin{figure}[h]
\centering
\includegraphics[width=\textwidth]{../pix/crlb_ratio.pdf}
\caption{Ratio of median CRLB (synchronized / asynchronous) as a
  function of connectivity $\theta$ and seasonal forcing amplitude
  $\delta$, for similar (left) and different (right) initial
  conditions. Darker shading indicates a larger ratio. The two panels
  look qualitatively similar, but the color scales differ by orders of
  magnitude: with similar ICs, synchronization inflates the CRLB far
  more than with different ICs.}
\label{fig:crlb_ratio}
\end{figure}

\subsection{Likelihood Surfaces: Simulated and Real Data}

Figure~\ref{fig:likelihood} shows log-likelihood surfaces for both
simulated (left) and real (right) data. In the simulated panel, we
use initial conditions drawn from the empirical distribution of
fitted values from the real-data analysis. Three scenarios are
compared: synchronized forcing with identical initial conditions
($\Delta\varphi = 0$), which produces a flat likelihood; a small
phase difference ($\Delta\varphi = \pi/45$), corresponding to the
median phase difference between US state pairs (excluding Alaska,
Hawaii, and Arizona, for which reliable absolute humidity data are
unavailable); and a large phase difference ($\Delta\varphi = \pi$),
which sharply peaks at the true $\theta$.

The right panel shows profile likelihood surfaces for real state
pairs. California and New York ($N_{\text{CA}} \approx 2
N_{\text{NY}}$) and Georgia and Ohio ($N_{\text{GA}} \approx
N_{\text{OH}}$) are chosen to illustrate the effect of population
size differences. In all cases, the likelihood is nearly flat with
respect to $\theta$, confirming the theoretical prediction that
synchronized temperate influenza epidemics carry insufficient
information for connectivity inference. However, our simulated
analysis reveals that symmetry breaking---whether through phase
differences, different initial conditions, or even different
population sizes---can restore identifiability
(SI~Fig.~S5). This suggests that for disease systems where
such asymmetries are naturally present, connectivity inference from
case data may be feasible.

\begin{figure}[h]
\centering
\includegraphics[width=\textwidth]{../pix/likelihood_surface.pdf}
\caption{Log-likelihood as a function of $\theta$. Left: simulated
  data with initial conditions drawn from the empirical distribution
  of real-data fits. Synchronized forcing with identical ICs produces
  a flat likelihood; increasing the phase difference ($\pi/45$
  corresponds to the median US state-pair difference) or using a large
  phase offset ($\pi$) progressively sharpens the peak around the
  true $\theta = 0.05$ (red dashed line). Right: profile likelihoods
  for real US state pairs. GA$\times$OH (similar populations) shows a
  flat surface; CA$\times$NY ($N_{\text{CA}} \approx 2 N_{\text{NY}}$)
  shows more curvature, suggesting that population asymmetry partially
  restores identifiability.}
\label{fig:likelihood}
\end{figure}

\subsection{US Influenza: Connectivity Cannot Be Inferred}

We estimated $\theta$ and its CRLB for all $\binom{47}{2} = 1081$
state pairs across 12 seasons. Figure~\ref{fig:theta_vs_crlb}~(left)
shows $\hat{\theta}$ against $\sqrt{\text{CRLB}}$, the theoretical
lower bound on the standard deviation of any unbiased estimator of
$\theta$. The dashed line marks the 5\% significance boundary: points
above it cannot reject $\theta = 0$. The vast majority of state pairs
fall in the non-significant region.

Figure~\ref{fig:theta_vs_crlb}~(right) shows the distribution of
$\hat\theta$ values across all state pairs. The estimates span the
full range $[0, 0.5]$ with an approximately uniform distribution,
consistent with the flat likelihood surfaces predicted by the
CRLB analysis: when the data contain no information about $\theta$,
the optimizer returns whatever value it happens to converge to, with
no tendency toward the truth.

\begin{figure}[h]
\centering
\includegraphics[width=\textwidth]{../pix/theta_vs_crlb.pdf}
\caption{Left: estimated connectivity $\hat{\theta}$ vs.\
  $\sqrt{\text{CRLB}}$ for all US state pairs. The dashed line is the
  significance boundary: points above it cannot reject $\theta = 0$
  at the 5\% level. Non-significant pairs are black circles;
  significant pairs are red stars. Right: distribution of
  $\hat\theta$ across all state pairs, showing an approximately
  uniform spread consistent with uninformative data.}
\label{fig:theta_vs_crlb}
\end{figure}

Representative model reconstructions for individual state pairs are
shown in SI~Fig.~S4.
